\chapter{समूह सिद्धांत के अनुप्रयोग}\label{ch:b3:group-theory-applied}

वायु के हर दस लाख अणुओं में कार्बन डाइऑक्साइड के कुछ सौ अणु होते हैं, और
नाइट्रोजन तथा ऑक्सीजन लगभग बाकी सब; फिर भी कार्बन डाइऑक्साइड ही, न कि
नाइट्रोजन या ऑक्सीजन, वह अवरक्त प्रकाश अवशोषित करती है जिसे गर्म भूमि
अंतरिक्ष की ओर भेजती है। कोई अणु अवरक्त प्रकाश केवल ऐसे कंपन द्वारा अवशोषित करता है जो
उसका द्विध्रुव आघूर्ण बदले, और कोई कंपन ऐसा करता है या नहीं, यह केवल
सममिति तय करती है: नाइट्रोजन का एकमात्र कंपन ऐसा नहीं कर सकता, कार्बन डाइऑक्साइड के
चार में से दो कर सकते हैं। यह अध्याय
\cref{ch:b3:point-groups} की \omterm{def:b3:point-groups:irrep}{अभिलक्षण सारणियों} को औज़ारों में बदलता है: यह निरूपणों का
लघुकरण करता है, सममिति-अनुकूलित कक्षक बनाता है, कंपनों को गिनता और अंकित करता है,
और अवरक्त तथा रामन स्पेक्ट्रमिकी के वरण नियम व्युत्पन्न करता है।

\begin{recall}
\Cref{ch:b3:point-groups} ने \omterm{def:b3:point-groups:operation}{सममिति संक्रियाएँ}, \omterm{def:b3:point-groups:point-group}{बिंदु समूह},
वर्ग, निरूपण, \omterm{def:b3:point-groups:representation}{अभिलक्षण} और \omterm{def:b3:point-groups:irrep}{अभिलक्षण सारणियाँ} परिभाषित कीं। द्वितीय वर्ष के
खंड ने \ce{H2O}, \ce{NH3} और \ce{CH4} के आण्विक कक्षक
अणुखंड कक्षकों और हाइड्रोजन $1s$
कक्षकों के सममिति-अनुकूलित संयोजनों से बनाए, और उन्हें ऐसे अंकनों से नाम दिया जो यहाँ व्युत्पन्न
किए गए हैं। प्रथम वर्ष के खंड ने तरंग-संख्याओं में अवरक्त स्पेक्ट्रम पढ़े और अणु की
ध्रुवणीयता परिभाषित की।
\end{recall}

\section{निरूपण का लघुकरण}

किसी अणु से जुड़े फलनों या सदिशों का कोई भी समुच्चय उसके \omterm{def:b3:point-groups:point-group}{बिंदु समूह} का एक
निरूपण वहन करता है। अधिकांश लघुकरणीय होते हैं।

\begin{definition}[लघुकरणीय निरूपण, पूर्णतः सममित निरूपण]\label{def:b3:group-theory-applied:reducible}
\emph{लघुकरणीय निरूपण}\index{लघुकरणीय निरूपण} वह है
जिसके सभी आव्यूहों को एक ही आधार-परिवर्तन से एक ही
खंड-विकर्ण रूप में लाया जा सके; तब वह \omterm{def:b3:point-groups:irrep}{अलघुकरणीय निरूपणों} का योग है,
जिसे $\Gamma = \sum_in_i\Gamma_i$ लिखा जाता है। वह \omterm{def:b3:point-groups:irrep}{अलघुकरणीय निरूपण} जिसके
सभी \omterm{def:b3:point-groups:representation}{अभिलक्षण} 1 हैं, \emph{पूर्णतः सममित
निरूपण}\index{पूर्णतः सममित निरूपण} है ($A_1$, $A_{1g}$,
$A_1'$\dots)।
\end{definition}

\begin{theorem}[महान लांबिकता प्रमेय]\label{thm:b3:group-theory-applied:great-orthogonality}
दो \omterm{def:b3:point-groups:irrep}{अलघुकरणीय निरूपणों} $\Gamma_i$, $\Gamma_j$ के लिए, जिनकी विमाएँ $d_i$,
$d_j$ हैं, जो कोटि $h$ वाले समूह के हैं और ऐकिक आव्यूहों में लिखे गए हैं,
\[
\sum_R D_i(R)_{mn}^*\,D_j(R)_{m'n'} = \frac{h}{d_i}\,\delta_{ij}\,\delta_{mm'}\,\delta_{nn'}.
\]
\end{theorem}
\admitted

उपपत्ति, शूर की प्रमेयिकाओं द्वारा, अधिक उन्नत पाठ्यक्रमों में दी जाती है;
\omterm{def:b3:point-groups:representation}{अभिलक्षणों} के लिए उसके परिणाम ही रसायन विज्ञान में प्रयुक्त होते हैं, और इस पुस्तक की
हर \omterm{def:b3:point-groups:irrep}{अभिलक्षण सारणी} उनके विरुद्ध जाँची गई है।

\begin{corollary}[अभिलक्षणों की लांबिकता]\label{cor:b3:group-theory-applied:characters}
$g_c$ को वर्ग $c$ की संक्रियाओं की संख्या लेने पर,
\[
\sum_cg_c\,\chi_i(c)^*\chi_j(c) = h\,\delta_{ij},
\]
और \omterm{def:b3:point-groups:irrep}{अलघुकरणीय निरूपणों} की संख्या वर्गों की संख्या के बराबर है,
और $\sum_id_i^2 = h$।
\end{corollary}

\begin{proof}
प्रमेय में $m = n$ और $m' = n'$ रखिए और $m$ तथा $m'$ पर योग कीजिए: बायाँ
पक्ष $\sum_R\chi_i(R)^*\chi_j(R)$ बन जाता है, दायाँ पक्ष $\frac{h}{d_i}\delta_{ij}
\sum_{m}\delta_{mm} = h\delta_{ij}$; संक्रियाओं को वर्ग के अनुसार समूहित कीजिए। इस प्रकार
सारणी की पंक्तियाँ एक ऐसे आकाश में लांबिक सदिश हैं जिसकी विमा
वर्गों की संख्या है, अतः \omterm{def:b3:point-groups:irrep}{अलघुकरणीय निरूपण} अधिक से अधिक उतने ही होते हैं
जितने वर्ग; समता और $\sum d_i^2 = h$ प्रमेय को
नियमित निरूपण पर लगाने से आते हैं (अभ्यास 10)।
\end{proof}

\begin{theorem}[लघुकरण सूत्र]\label{thm:b3:group-theory-applied:reduction}
निरूपण $\Gamma$, जिसके \omterm{def:b3:point-groups:representation}{अभिलक्षण} $\chi(c)$ हैं, \omterm{def:b3:point-groups:irrep}{अलघुकरणीय
निरूपण} $\Gamma_i$ को ठीक इतनी बार समाहित करता है:
\[
n_i = \frac1h\sum_cg_c\,\chi(c)\,\chi_i(c)^*
\]
यही \emph{लघुकरण सूत्र}\index{लघुकरण सूत्र} है।
\end{theorem}

\begin{proof}
चूँकि $\Gamma = \sum_jn_j\Gamma_j$, उसके \omterm{def:b3:point-groups:representation}{अभिलक्षण} $\chi(c) = \sum_jn_j\chi_j(c)$ हैं।
$g_c\chi_i(c)^*$ से गुणा कीजिए, वर्गों पर योग कीजिए, और उपप्रमेय का उपयोग कीजिए:
$\sum_cg_c\chi(c)\chi_i(c)^* = \sum_jn_jh\delta_{ij} = hn_i$।
\end{proof}

\begin{example}[अमोनिया के हाइड्रोजन कक्षक]\label{ex:b3:group-theory-applied:nh3}
\ce{NH3} के हाइड्रोजनों के तीन $1s$ कक्षकों को आधार लीजिए। $E$
तीनों को उनके स्थान पर छोड़ता है ($\chi = 3$), $C_3$ तीनों को हिलाता है ($\chi = 0$), कोई
$\sigma_v$ एक को स्थान पर छोड़ता है और दो की अदला-बदली करता है ($\chi = 1$)। $C_{3v}$ सारणी से:
$n_{A_1} = \frac16(3 + 0 + 3\times1) = 1$, $n_{A_2} = \frac16(3 + 0 - 3) = 0$,
$n_E = \frac16(6 + 0 + 0) = 1$। अतः $\Gamma_{\mathrm H} = A_1 + E$: तीनों
हाइड्रोजन एक पूर्णतः सममित संयोजन और एक अपभ्रष्ट युग्म देते हैं।
\end{example}

\begin{method}[निरूपण का लघुकरण]\label{met:b3:group-theory-applied:reduce}
\begin{enumerate}
  \item हर वर्ग के लिए \omterm{def:b3:point-groups:representation}{अभिलक्षण} ज्ञात कीजिए: परमाणु-केंद्रित फलनों के समुच्चय के लिए
        यह स्थान पर छोड़े गए फलनों की संख्या है, हर एक को उस चिह्न के साथ गिनकर
        जो उसे मिलता है।
  \item हर \omterm{def:b3:point-groups:irrep}{अलघुकरणीय निरूपण} की पंक्ति के साथ \omterm{thm:b3:group-theory-applied:reduction}{लघुकरण सूत्र} लगाइए,
        हर वर्ग को उसके आकार से भारित करते हुए।
  \item जाँचिए: परिणाम अऋणात्मक पूर्णांक हों, और $\sum_in_id_i$
        $E$ के अंतर्गत \omterm{def:b3:point-groups:representation}{अभिलक्षण} (आधार के आकार) के बराबर हो।
\end{enumerate}
\end{method}

\section{सममिति-अनुकूलित कक्षक}

\begin{definition}[प्रक्षेप संकारक]\label{def:b3:group-theory-applied:projection}
\omterm{def:b3:point-groups:irrep}{अलघुकरणीय निरूपण} $\Gamma_i$ पर \emph{प्रक्षेप संकारक}\index{प्रक्षेप संकारक}
यह है:
\[
\hat P_i = \frac{d_i}{h}\sum_R\chi_i(R)^*\,\hat R,
\]
जहाँ $\hat R$ किसी फलन पर संक्रिया $R$ लगाता है।
\end{definition}

\begin{proposition}[प्रक्षेप संकारक क्या करता है]\label{prop:b3:group-theory-applied:projection}
किसी भी फलन $f$ पर लगाने पर $\hat P_if$ या तो शून्य है या ऐसा फलन जो
$\Gamma_i$ की तरह रूपांतरित होता है (एक सममिति-अनुकूलित संयोजन); एकविमीय
$\Gamma_i$ के लिए वह, चिह्न $\chi_i(S)$ तक, हर
संक्रिया $S$ से अपरिवर्तित रहता है।
\end{proposition}

\begin{proof}[आंशिक उपपत्ति]
$d_i = 1$ के लिए: $\hat S\hat P_if = \frac1h\sum_R\chi_i(R)\hat S\hat Rf$। $R' = SR$ रखिए,
जो $R$ की तरह पूरे समूह पर चलता है; $\chi_i(R) = \chi_i(S^{-1}R') =
\chi_i(S)\chi_i(R')$, एकविमीय निरूपण के लिए (\omterm{def:b3:point-groups:representation}{अभिलक्षण} ही
आव्यूह हैं)। अतः $\hat S\hat P_if = \chi_i(S)\hat P_if$। सामान्य
स्थिति महान लांबिकता प्रमेय का उपयोग करती है और स्वीकार की गई है।
\end{proof}

\begin{method}[सममिति-अनुकूलित संयोजन बनाना]\label{met:b3:group-theory-applied:salc}
\begin{enumerate}
  \item समुच्चय का एक फलन, $h_1$, चुनिए और सारणीबद्ध कीजिए कि हर
        संक्रिया उसे किस फलन पर भेजती है।
  \item हर \omterm{def:b3:point-groups:irrep}{अलघुकरणीय निरूपण} के लिए, हर प्रतिबिंब को संक्रिया के
        \omterm{def:b3:point-groups:representation}{अभिलक्षण} से गुणा कीजिए, जोड़िए, और प्रसामान्यीकृत कीजिए (परमाणु फलनों के बीच
        अतिव्यापन की उपेक्षा करते हुए)।
  \item अपभ्रष्ट निरूपण के लिए, दूसरा सदस्य पाने हेतु दूसरे फलन का प्रक्षेप
        कीजिए, फिर लांबिकीकरण कीजिए।
\end{enumerate}
\end{method}

\begin{example}[अमोनिया और जल के संयोजन]\label{ex:b3:group-theory-applied:salcs}
\ce{NH3} में $\hat P_{A_1}h_1 \propto h_1 + h_2 + h_3$ (सभी \omterm{def:b3:point-groups:representation}{अभिलक्षण} 1)। $E$ के लिए
(\omterm{def:b3:point-groups:representation}{अभिलक्षण} $2, -1, 0$), $\hat P_Eh_1 \propto 2h_1 - h_2 - h_3$; $h_2$ का प्रक्षेप करके
लांबिकीकरण करने पर साथी $h_2 - h_3$ मिलता है। प्रसामान्यीकृत: $a_1 =
\frac{1}{\sqrt3}(h_1 + h_2 + h_3)$, $e = \frac{1}{\sqrt6}(2h_1 - h_2 - h_3)$,
$\frac{1}{\sqrt2}(h_2 - h_3)$। \ce{H2O} में, जिसे $xz$ तल में रखा गया है, $h_1 + h_2$
$a_1$ है और $h_1 - h_2$, $b_1$ है (वह $C_2$ के अंतर्गत और $\sigma_v'(yz)$ के अंतर्गत चिह्न बदलता है,
जो हाइड्रोजनों की अदला-बदली करते हैं)। अणु को $yz$ तल में रखने पर, जैसा कुछ
पुस्तकें पसंद करती हैं, अंकन $b_1$ और $b_2$ सर्वत्र परस्पर बदल जाते हैं।
\end{example}

\begin{omfigure}
\begin{tikzpicture}[font=\small]
  \foreach \x/\lab/\ca/\cb/\cc in {0/{$a_1$}/1/1/1, 4/{$e$}/2/-1/-1, 8/{$e$}/0/1/-1}{
    \begin{scope}[xshift=\x cm]
      \foreach \a/\c in {90/\ca,210/\cb,330/\cc}{
        \pgfmathsetmacro{\r}{0.18+0.13*abs(\c)}
        \ifdim \c pt > 0pt \fill[omDef!55, draw=omDef] (\a:1.0) circle (\r);
        \else \ifdim \c pt < 0pt \fill[omThm!25, draw=omThm] (\a:1.0) circle (\r);
        \else \draw[gray, densely dotted] (\a:1.0) circle (0.12); \fi\fi
      }
      \node[aN, minimum size=4.6mm] at (0,0) {};
      \node[anchor=north] at (0,-1.45) {\lab};
    \end{scope}}
  \node[anchor=north, font=\footnotesize] at (0,-1.85) {$h_1 + h_2 + h_3$};
  \node[anchor=north, font=\footnotesize] at (4,-1.85) {$2h_1 - h_2 - h_3$};
  \node[anchor=north, font=\footnotesize] at (8,-1.85) {$h_2 - h_3$};
\end{tikzpicture}

{\small अमोनिया के तीन हाइड्रोजन $1s$ कक्षकों के सममिति-अनुकूलित संयोजन,
$C_3$ अक्ष के अनुदिश देखे गए (ऊपर $h_1$)। नीला: धनात्मक
गुणांक, लाल: ऋणात्मक, हर चकती का आकार
गुणांक के साथ बढ़ता है; बिंदुकित: शून्य।}
\end{omfigure}

\begin{example}[किसी धातु के चारों ओर छह लिगैंड]\label{ex:b3:group-theory-applied:ml6}
अष्टफलक के शीर्षों से किसी धातु की ओर इंगित छह लिगैंड $\sigma$ कक्षकों के लिए,
$O_h$ के वर्गों ($E$, $8C_3$, $6C_2$, $6C_4$,
$3C_2$, $i$, $6S_4$, $8S_6$, $3\sigma_h$, $6\sigma_d$) पर \omterm{def:b3:point-groups:representation}{अभिलक्षण} $6, 0, 0, 2, 2, 0, 0, 0, 4,
2$ हैं, और लघुकरण $\Gamma_\sigma = A_{1g} + E_g + T_{1u}$ देता है। धातु के $s$
($a_{1g}$), $d_{z^2}$ और $d_{x^2-y^2}$ ($e_g$) तथा $p$ ($t_{1u}$) कक्षकों को लिगैंड
साथी मिलते हैं; उसके $d_{xy}$, $d_{xz}$, $d_{yz}$ ($t_{2g}$) को कोई नहीं मिलता और वे
$\sigma$ में अनाबंधी रहते हैं --- यही द्वितीय वर्ष के खंड के $t_{2g}$/$e_g$ विपाटन का मूल है,
जिसे \cref{ch:b3:organometallic-bonding,ch:b3:complex-spectra-magnetism} में विकसित किया गया है।
\end{example}

\begin{omfigure}
\begin{tikzpicture}[font=\small, yscale=0.95]
  % O orbitals (left), MOs (centre), H combinations (right); schematic energies
  \foreach \y/\l in {0/{$2s\ (a_1)$}}{\draw[omstate] (0,\y) -- (1.2,\y); \node[anchor=east] at (-0.05,\y) {\l};}
  \draw[omstate] (0,3.0) -- (0.36,3.0) (0.42,3.0) -- (0.78,3.0) (0.84,3.0) -- (1.2,3.0);
  \node[anchor=east] at (-0.05,3.0) {$2p\ (b_1, a_1, b_2)$};
  \draw[omstate] (6.8,3.9) -- (8.0,3.9); \node[anchor=west] at (8.05,3.9) {$h_1 - h_2\ (b_1)$};
  \draw[omstate] (6.8,3.6) -- (8.0,3.6); \node[anchor=west] at (8.05,3.5) {$h_1 + h_2\ (a_1)$};
  \foreach \y/\l in {-0.6/{$2a_1$}, 1.6/{$1b_1$}, 2.4/{$3a_1$}, 3.0/{$1b_2$}, 5.6/{$4a_1$}, 6.2/{$2b_1$}}{
    \draw[omstate] (3.4,\y) -- (4.6,\y); \node[anchor=west] at (4.65,\y) {\l};}
  \foreach \y in {-0.6,1.6,2.4,3.0}{\node[font=\scriptsize] at (4.0,\y+0.16) {$\uparrow\downarrow$};}
  \draw[densely dotted, gray] (1.2,0) -- (3.4,-0.6) (1.2,3.0) -- (3.4,1.6) (1.2,3.0) -- (3.4,2.4) (1.2,3.0) -- (3.4,3.0)
     (6.8,3.6) -- (4.6,-0.6) (6.8,3.9) -- (4.6,1.6) (6.8,3.6) -- (4.6,2.4) (1.2,3.0) -- (3.4,6.2) (6.8,3.9) -- (4.6,6.2) (6.8,3.6) -- (4.6,5.6);
  \node[anchor=north] at (0.6,-1.0) {O};
  \node[anchor=north] at (4.0,-1.0) {\ce{H2O}};
  \node[anchor=north] at (7.4,-1.0) {2 H};
  \draw[->] (-2.3,-0.9) -- (-2.3,6.4) node[above] {$E$};
\end{tikzpicture}

{\small जल के संयोजकता आण्विक कक्षक (योजनात्मक ऊर्जाएँ, अणु
$xz$ तल में)। केवल समान सममिति के कक्षक मिश्रित होते हैं: $b_1$ हाइड्रोजन
संयोजन $2p_x$ के साथ, $a_1$ वाला $2s$ और $2p_z$ के साथ; $2p_y$ ($b_2$) का कोई
साथी नहीं है और वह एक अनाबंधी एकाकी युग्म, उच्चतम अधिकृत कक्षक, बना रहता है।}
\end{omfigure}

\section{कंपनिक विधाएँ}

\begin{definition}[प्रसामान्य विधा]\label{def:b3:group-theory-applied:normal-mode}
किसी अणु की \emph{प्रसामान्य विधा}\index{प्रसामान्य विधा} एक सामूहिक
कंपन है जिसमें सभी परमाणु एक ही आवृत्ति पर और एक ही कला में दोलन करते हैं,
द्रव्यमान-भारित हेसियन के अभिलक्षणिक सदिश के अनुदिश
(\cref{prop:b3:computational-chemistry:hessian})। हर प्रसामान्य विधा
\omterm{def:b3:point-groups:point-group}{बिंदु समूह} के किसी \omterm{def:b3:point-groups:irrep}{अलघुकरणीय निरूपण} की तरह रूपांतरित होती है।
\end{definition}

\begin{proposition}[सभी गतियों का निरूपण]\label{prop:b3:group-theory-applied:gamma-3n}
हर परमाणु से तीन विस्थापन सदिश $x$, $y$, $z$ जोड़िए। किसी संक्रिया
$R$ के अंतर्गत इस $3N$-विमीय निरूपण का \omterm{def:b3:point-groups:representation}{अभिलक्षण} यह है:
\[
\chi_{3N}(R) = (\text{स्थान पर छोड़े गए परमाणुओं की संख्या}) \times \chi_{xyz}(R),
\]
जहाँ $\chi_{xyz} = 1 + 2\cos\theta$, $\theta$ के घूर्णन के लिए, और $-1 + 2\cos\theta$
अनुचित घूर्णन के लिए (अतः $E$ के लिए 3, $-1$, $C_2$ के लिए, 0, $C_3$ के लिए, 1, $\sigma$ के लिए,
$-3$, $i$ के लिए)।
\end{proposition}

\begin{proof}
किसी अन्य स्थिति पर ले जाया गया परमाणु अपने सदिशों को आव्यूह के विकर्ण से बाहर ले
जाता है: अनुरेख में उसका कोई योगदान नहीं। स्थान पर छोड़े गए परमाणु के तीन सदिश
$3\times3$ आव्यूह द्वारा रूपांतरित होते हैं, जो $R$ का आव्यूह है और जिसका अनुरेख
अक्ष के अनुदिश $z$ वाले निर्देश-तंत्र में परिकलित होता है: घूर्णन के लिए $\cos\theta + \cos\theta + 1$,
अनुचित घूर्णन के लिए अंतिम प्रविष्टि $-1$ हो जाती है।
\end{proof}

\begin{proposition}[कंपनों की गिनती]\label{prop:b3:group-theory-applied:mode-count}
$\Gamma_{\mathrm{vib}} = \Gamma_{3N} - \Gamma_{\mathrm{trans}} - \Gamma_{\mathrm{rot}}$, जहाँ
$\Gamma_{\mathrm{trans}}$, $(x, y, z)$ का निरूपण है और $\Gamma_{\mathrm{rot}}$,
$(R_x, R_y, R_z)$ का; अरैखिक अणु के $3N - 6$ कंपन होते हैं, रैखिक
के $3N - 5$।
\end{proposition}

\begin{proof}
$3N$ विस्थापन हर गति का वर्णन करते हैं: द्रव्यमान केंद्र के तीन स्थानांतरण,
तीन घूर्णन (रैखिक अणु के लिए दो, क्योंकि अपने ही अक्ष के परितः घूमने से
कोई नाभिक नहीं हिलता), और कंपन। संक्रियाओं के अंतर्गत ये उपसमष्टियाँ
मिश्रित नहीं होतीं, अतः उनके निरूपण जुड़ते हैं।
\end{proof}

\begin{example}[जल]\label{ex:b3:group-theory-applied:water}
% ledger: vib:H2O.sym, vib:H2O.bend, vib:H2O.asym
$C_{2v}$ में \ce{H2O}, अणु $xz$ तल में। स्थान पर छोड़े गए परमाणु: 3, 1, 3, 1;
$\chi_{xyz}$: $3, -1, 1, 1$; अतः $\chi_{3N} = 9, -1, 3, 1$। लघुकरण $3A_1 +
A_2 + 3B_1 + 2B_2$ देता है। स्थानांतरण $A_1 + B_1 + B_2$ और घूर्णन $A_2 + B_1 + B_2$
घटाने पर: $\Gamma_{\mathrm{vib}} = 2A_1 + B_1$। दोनों $a_1$ विधाएँ सममित
तनन (\qty{3657}{cm^{-1}}) और बंकन (\qty{1595}{cm^{-1}}) हैं; $b_1$ विधा
प्रतिसममित तनन (\qty{3756}{cm^{-1}}) है।
\end{example}

\begin{omfigure}
\begin{tikzpicture}[font=\small, >=Stealth]
  \foreach \x/\name/\lab in {0/sym/{$\nu_1$, $a_1$, सममित तनन}, 4.3/bend/{$\nu_2$, $a_1$, बंकन}, 8.6/asym/{$\nu_3$, $b_1$, प्रतिसममित तनन}}{
    \begin{scope}[xshift=\x cm]
      \node[aO] (o\name) at (0,0.35) {};
      \node[aH] (a\name) at (-0.8,-0.25) {};
      \node[aH] (b\name) at (0.8,-0.25) {};
      \begin{scope}[on background layer]\draw[bond] (o\name.center) -- (a\name.center) (o\name.center) -- (b\name.center);\end{scope}
      \node[anchor=north, align=center, font=\footnotesize] at (0,-0.9) {\lab};
    \end{scope}}
  % symmetric stretch: H outward along the bonds, O up a little
  \draw[->, omThm, thick] (-0.95,-0.36) -- (-1.42,-0.71);
  \draw[->, omThm, thick] (0.95,-0.36) -- (1.42,-0.71);
  \draw[->, omThm, thick] (0,0.74) -- (0,0.98);
  % bend: each H moves perpendicular to its bond, closing the angle; O up
  \draw[->, omThm, thick] (4.3-0.92,-0.36) -- (4.3-0.62,-0.76);
  \draw[->, omThm, thick] (4.3+0.92,-0.36) -- (4.3+0.62,-0.76);
  \draw[->, omThm, thick] (4.3,0.74) -- (4.3,0.98);
  % antisymmetric stretch: one H out, the other in, O sideways
  \draw[->, omThm, thick] (8.6-0.95,-0.36) -- (8.6-1.42,-0.71);
  \draw[->, omThm, thick] (8.6+1.02,0.0) -- (8.6+0.70,0.24);
  \draw[->, omThm, thick] (8.6-0.05,0.80) -- (8.6+0.25,0.80);
\end{tikzpicture}

{\small जल की तीन \omterm{def:b3:group-theory-applied:normal-mode}{प्रसामान्य विधाएँ} (तीर: विस्थापन, पैमाने पर
नहीं)। दोनों $a_1$ विधाएँ अणु की पूरी सममिति बनाए रखती हैं; $b_1$
विधा $C_2$ के अंतर्गत चिह्न बदलती है।}
\end{omfigure}

\begin{example}[चार और अणु]\label{ex:b3:group-theory-applied:table}
यही प्रक्रिया (इस अध्याय के चित्र-आँकड़ों में परिकलित और जाँची गई)
यह देती है:
\begin{center}\small
\begin{tabular}{llll}
\hline
अणु & समूह & $\Gamma_{\mathrm{vib}}$ & विधाएँ \\ \hline
\ce{NH3} & $C_{3v}$ & $2A_1 + 2E$ & 6 \\
\ce{CH4} & $T_d$ & $A_1 + E + 2T_2$ & 9 \\
\ce{CO2} & $D_{\infty h}$ ($D_{2h}$ के माध्यम से) & $\Sigma_g^+ + \Sigma_u^+ + \Pi_u$ ($A_g + B_{1u} + B_{2u} + B_{3u}$) & 4 \\
\ce{BF3} & $D_{3h}$ & $A_1' + 2E' + A_2''$ & 6 \\
\ce{XeF4} & $D_{4h}$ & $A_{1g} + B_{1g} + B_{2g} + A_{2u} + B_{2u} + 2E_u$ & 9 \\
\ce{SF6} & $O_h$ & $A_{1g} + E_g + T_{2g} + 2T_{1u} + T_{2u}$ & 15 \\ \hline
\end{tabular}
\end{center}
रैखिक अणु के लिए अनंत समूह के स्थान पर उसका \omterm{def:b3:point-groups:class}{उपसमूह} $D_{2h}$ लिया जाता है,
जिसमें अपभ्रष्ट $\Pi_u$ बंकन $B_{2u} + B_{3u}$ में विभक्त हो जाता है।
\end{example}

\section{वरण नियम}

\begin{definition}[प्रत्यक्ष गुणनफल]\label{def:b3:group-theory-applied:direct-product}
दो निरूपणों का \emph{प्रत्यक्ष गुणनफल}\index{प्रत्यक्ष गुणनफल} $\Gamma_i\otimes\Gamma_j$
वह निरूपण है जिसे उनके आधार फलनों के गुणनफल वहन करते हैं; उसके \omterm{def:b3:point-groups:representation}{अभिलक्षण}
गुणनफल $\chi_i(R)\chi_j(R)$ हैं।
\end{definition}

\begin{theorem}[लुप्त होने वाले समाकल]\label{thm:b3:group-theory-applied:vanishing-integral}
पूरे आकाश पर समाकल $\int f_1f_2f_3\,\dd\tau$ शून्येतर केवल तभी हो सकता है जब
\omterm{def:b3:group-theory-applied:direct-product}{प्रत्यक्ष गुणनफल} $\Gamma_1\otimes\Gamma_2\otimes\Gamma_3$ में \omterm{def:b3:group-theory-applied:reducible}{पूर्णतः
सममित निरूपण} समाहित हो।
\end{theorem}

\begin{proof}
\omterm{def:b3:point-groups:operation}{सममिति संक्रिया} केवल आकाश के बिंदुओं को पुनः अंकित करती है, अतः समाकल्य को किसी भी
$R$ से रूपांतरित करने पर समाकल का मान अपरिवर्तित रहता है। रूपांतरित समाकल्य का
समूह पर औसत लीजिए: हर \omterm{def:b3:point-groups:irrep}{अलघुकरणीय निरूपण} से संबंधित समाकल्य के घटकों का औसत
$\frac1h\sum_R\hat R(\cdot)$ है, जो \omterm{def:b3:group-theory-applied:reducible}{पूर्णतः सममित
निरूपण} पर प्रक्षेप है (\cref{def:b3:group-theory-applied:projection}, जहाँ सभी
$\chi = 1$)। हर अन्य घटक का औसत शून्य होता है; यदि गुणनफल में कोई
पूर्णतः सममित भाग नहीं है, तो समाकल शून्य है।
\end{proof}

\begin{definition}[अवरक्त-सक्रिय, रामन-सक्रिय]\label{def:b3:group-theory-applied:activity}
कोई मूल कंपन \emph{अवरक्त-सक्रिय}\index{अवरक्त-सक्रिय} है यदि वह
अवरक्त प्रकाश अवशोषित कर सके, \emph{रामन-सक्रिय}\index{रामन-सक्रिय} यदि वह
रामन स्पेक्ट्रम में प्रकट हो (\cref{ch:b3:rovibrational-spectroscopy})।
\end{definition}

\begin{corollary}[अवरक्त और रामन वरण नियम]\label{cor:b3:group-theory-applied:ir-raman}
कोई मूल संक्रमण (मूल स्तर से, जो पूर्णतः सममित है, सममिति $\Gamma$ वाली किसी
विधा के एक क्वांटम तक) \omterm{def:b3:group-theory-applied:activity}{अवरक्त-सक्रिय} केवल तभी है जब $\Gamma$,
$x$, $y$ या $z$ का निरूपण हो, और \omterm{def:b3:group-theory-applied:activity}{रामन-सक्रिय} केवल तभी जब $\Gamma$ किसी
द्विघाती फलन ($x^2$, $xy$, \dots) का निरूपण हो।
\end{corollary}

\begin{proof}
अवशोषण की तीव्रता में $\int\psi_0\,\mu_k\,\psi_1\dd\tau$ आता है, जहाँ द्विध्रुव
घटक $\mu_k$, $x$, $y$, $z$ की तरह रूपांतरित होते हैं; $\psi_0$ पूर्णतः सममित है
और $\psi_1$, $\Gamma$ की तरह रूपांतरित होता है। प्रमेय के अनुसार समाकल शून्येतर
केवल तभी हो सकता है जब $\Gamma\otimes\Gamma_{\mu_k}$ में $A_1$ हो, जो केवल तब होता है जब $\Gamma =
\Gamma_{\mu_k}$ (दो \omterm{def:b3:point-groups:irrep}{अलघुकरणीय निरूपणों} के गुणनफल में \omterm{def:b3:group-theory-applied:reducible}{पूर्णतः सममित
निरूपण} केवल तभी होता है जब वे बराबर हों, वास्तविक \omterm{def:b3:point-groups:representation}{अभिलक्षणों} के लिए)। रामन
प्रकीर्णन में ध्रुवणीयता के घटक $\alpha_{kl}$ आते हैं, जो
गुणनफलों $kl$ की तरह रूपांतरित होते हैं।
\end{proof}

\begin{proposition}[परस्पर अपवर्जन नियम]\label{prop:b3:group-theory-applied:mutual-exclusion}
\omterm{def:b3:point-groups:inversion}{प्रतिलोमन केंद्र} वाले अणु में कोई भी मूल कंपन \omterm{def:b3:group-theory-applied:activity}{अवरक्त-सक्रिय} और \omterm{def:b3:group-theory-applied:activity}{रामन-सक्रिय}
दोनों नहीं होता: यही \emph{परस्पर अपवर्जन नियम}\index{परस्पर अपवर्जन नियम} है।
\end{proposition}

\begin{proof}
\omterm{def:b3:point-groups:inversion}{प्रतिलोमन केंद्र} होने पर हर \omterm{def:b3:point-groups:irrep}{अलघुकरणीय निरूपण} $g$ या $u$ होता है।
निर्देशांक $x$, $y$, $z$, $i$ के अंतर्गत चिह्न बदलते हैं (वे $u$ हैं); द्विघाती
फलन नहीं बदलते ($g$)। कोई विधा \omterm{def:b3:group-theory-applied:activity}{अवरक्त-सक्रिय} केवल तभी है जब $u$ हो, \omterm{def:b3:group-theory-applied:activity}{रामन-सक्रिय} केवल तभी जब
$g$ हो।
\end{proof}

\begin{example}[कार्बन डाइऑक्साइड और ग्रीनहाउस प्रभाव]\label{ex:b3:group-theory-applied:co2}
% ledger: vib:CO2.sym, vib:CO2.bend, vib:CO2.asym, ir:CO2.asym.int, ir:CO2.bend.int
\ce{CO2} में \omterm{def:b3:point-groups:inversion}{प्रतिलोमन केंद्र} है। उसका सममित तनन ($\Sigma_g^+$,
\qty{1333}{cm^{-1}}) केवल \omterm{def:b3:group-theory-applied:activity}{रामन-सक्रिय} है; प्रतिसममित तनन ($\Sigma_u^+$,
\qty{2349}{cm^{-1}}) और बंकन ($\Pi_u$, \qty{667}{cm^{-1}}) केवल \omterm{def:b3:group-theory-applied:activity}{अवरक्त-सक्रिय} हैं।
बंकन \qty{15}{\micro m} के निकट अवशोषण करता है, पृथ्वी के
तापीय उत्सर्जन के शिखर के पास: यही कारण है कि कार्बन डाइऑक्साइड ग्रीनहाउस गैस है। \ce{N2} और
\ce{O2} का एकमात्र, $\Sigma_g^+$, कंपन है: अवरक्त-निष्क्रिय।
\end{example}

\begin{omfigure}
\begin{tikzpicture}
\begin{axis}[omir, width=0.92\linewidth, height=4.4cm, xmin=400, xmax=2600,
  ymin=0, ymax=1.3, ytick=\empty, xlabel={तरंग-संख्या / cm$^{-1}$},
  ylabel={संकेत}, legend cell align=left,
  legend style={font=\small, at={(0.98,0.95)}, anchor=north east, draw=none}]
\addplot[omDef, thick] table[x=nu, y=ir] {figdata/out/bachelor-3-co2-spectra.dat};
\addlegendentry{अवरक्त अवशोषण}
\addplot[omThm, thick, dashed] table[x=nu, y=raman] {figdata/out/bachelor-3-co2-spectra.dat};
\addlegendentry{रामन प्रकीर्णन}
\node[font=\small, anchor=south] at (axis cs:2349,1.02) {$\Sigma_u^+$};
\node[font=\small, anchor=south] at (axis cs:667,0.1) {$\Pi_u$};
\node[font=\small, anchor=south, omThm] at (axis cs:1333,1.02) {$\Sigma_g^+$};
\end{axis}
\end{tikzpicture}

{\small गैसीय \ce{CO2} के संश्लिष्ट स्पेक्ट्रम (बैंडों की स्थितियाँ और आपेक्षिक
अवरक्त तीव्रताएँ मापे गए मानों से; आकृतियाँ योजनात्मक, बिना
घूर्णी संरचना के)। परस्पर अपवर्जन: कोई बैंड दोनों में प्रकट नहीं होता।}
\end{omfigure}

\begin{method}[कंपनिक स्पेक्ट्रम की भविष्यवाणी]\label{met:b3:group-theory-applied:vibrations}
\begin{enumerate}
  \item \omterm{def:b3:point-groups:point-group}{बिंदु समूह} ज्ञात कीजिए; स्थान पर छोड़े गए परमाणुओं से $\chi_{3N}$ परिकलित कीजिए।
  \item लघुकरण कीजिए, स्थानांतरण और घूर्णन घटाइए: $\Gamma_{\mathrm{vib}}$।
  \item सारणी के दाहिने स्तंभों से हर विधा को \omterm{def:b3:group-theory-applied:activity}{अवरक्त-सक्रिय}
        ($x$, $y$, $z$) और/या \omterm{def:b3:group-theory-applied:activity}{रामन-सक्रिय} (द्विघाती) चिह्नित कीजिए; जो विधा दोनों में से कोई
        नहीं है वह मूक है।
  \item बैंड गिनिए: हर सक्रिय विधा का एक, अपभ्रष्ट विधा एक ही बैंड देती है।
\end{enumerate}
\end{method}

\begin{method}[CO तननों की गिनती]\label{met:b3:group-theory-applied:co-count}
किसी धातु कार्बोनिल में $n$ C--O आबंध सदिशों को आधार लीजिए (स्थान पर छोड़ा गया
परमाणु 1 गिना जाता है, बाकी कुछ नहीं): लघुकरण करके $\Gamma_{\mathrm{CO}}$ प्राप्त कीजिए, फिर
अवरक्त और रामन नियम लगाइए।
\qtyrange{1850}{2150}{cm^{-1}} के निकट प्रबल बैंडों की संख्या समावयवी की पहचान करती है।
\end{method}

\begin{example}[सिस और ट्रांस]\label{ex:b3:group-theory-applied:cis-trans}
cis-\ce{ML4(CO)2}, $C_{2v}$: \omterm{def:b3:point-groups:representation}{अभिलक्षण} $2, 0, 2, 0$, $\Gamma_{\mathrm{CO}} = A_1 + B_1$,
दो अवरक्त बैंड। trans-\ce{ML4(CO)2}, $D_{4h}$: $\Gamma_{\mathrm{CO}} = A_{1g} + A_{2u}$, एक
अवरक्त बैंड ($A_{2u}$) और एक रामन बैंड ($A_{1g}$)। अवरक्त में एक ही CO बैंड
का अर्थ है trans।
\end{example}

\section{इलेक्ट्रॉनिक संक्रमणों पर अनुप्रयोग}

यही प्रमेय इलेक्ट्रॉनिक संक्रमणों पर लागू होता है: अवस्थाओं $\Psi_1$ और $\Psi_2$ के बीच
इलेक्ट्रॉनिक संक्रमण केवल तभी अनुमत है जब $\Gamma_1\otimes\Gamma_2$ में
$x$, $y$ या $z$ का निरूपण हो, और उसका प्रकाश
उसी अक्ष के अनुदिश ध्रुवित होता है।

\begin{example}[जल और फ़ॉर्मेल्डिहाइड के संक्रमण]\label{ex:b3:group-theory-applied:electronic}
जल ($C_{2v}$, $xz$ तल) में $1b_2$ (एकाकी युग्म) से
$4a_1$ में एक इलेक्ट्रॉन का उन्नयन सममिति $B_2\otimes A_1 = B_2$ की उत्तेजित अवस्था देता है, जो $y$ है:
अनुमत, आण्विक तल के लंबवत ध्रुवित। फ़ॉर्मेल्डिहाइड
($C_{2v}$) में $n\to\pi^*$ संक्रमण ऑक्सीजन के तल-स्थित एकाकी युग्म ($b_1$) से
$\pi^*$ कक्षक तक जाता है, जो तल के लंबवत $p$ कक्षकों से बना है
($b_2$; अणु $xz$ तल में, C=O $z$ के अनुदिश): $B_1\otimes B_2 = A_2$,
जो $x$, $y$, $z$ में से कोई नहीं है: सममिति-वर्जित, इसीलिए उसका बैंड
(\cref{ch:b3:electronic-spectroscopy}) इतना दुर्बल है।
\end{example}

\begin{history}[विग्नर, बेथे और अणुओं की सममिति]
समूह सिद्धांत ने क्वांटम यांत्रिकी में यूजीन विग्नर के परमाण्वीय
स्पेक्ट्रमों पर कार्य (1927--31) और क्रिस्टलों में परमाण्वीय पदों के विपाटन पर हान्स
बेथे के शोधपत्र (1929) के साथ प्रवेश किया। ई॰ ब्राइट विल्सन ने 1934 में इसे आण्विक कंपनों पर
लागू किया; रॉबर्ट मुलिकेन ने कक्षकों और अवस्थाओं के वे अंकन दिए जो आज भी प्रयुक्त होते हैं।
\end{history}

\begin{inthelab}[अवरक्त और रामन साथ-साथ]
अवरक्त स्पेक्ट्रममापी वह प्रकाश मापता है जिसे नमूना अवशोषित करता है; रामन
स्पेक्ट्रममापी उसे लेज़र से प्रदीप्त करता है और विस्थापित तरंगदैर्घ्यों पर
प्रकीर्णित दुर्बल प्रकाश का विश्लेषण करता है। एक ही यौगिक पर दोनों अभिलिखित करना
सममिति केंद्र का चिरप्रचलित परीक्षण है: यदि कोई बैंड संपाती नहीं होता, तो अणु
संभवतः केंद्रसममित है। रामन लेज़र वर्ग 3B या 4 के होते हैं: किरण-पुंज
घेरकर रखा जाता है और संचालक उसके तरंगदैर्घ्य के लिए निर्धारित चश्मे पहनते हैं।
\end{inthelab}

\section{अभ्यास}

\begin{exercise}[$\star$]\label{exo:b3:group-theory-applied:1}
$C_{2v}$ के उस निरूपण का लघुकरण कीजिए जिसके \omterm{def:b3:point-groups:representation}{अभिलक्षण} $(5, -1, 3, 1)$ हैं, $(E, C_2,
\sigma_v(xz), \sigma_v'(yz))$ पर। विमा जाँचिए।
\end{exercise}

\begin{exercise}[$\star$]\label{exo:b3:group-theory-applied:2}
\ce{SO2}, जल की तरह, बंकित है। $\Gamma_{\mathrm{vib}}$ दीजिए और बताइए कि कौन-सी विधाएँ अवरक्त-
और \omterm{def:b3:group-theory-applied:activity}{रामन-सक्रिय} हैं।
\end{exercise}

\begin{exercise}[$\star$]\label{exo:b3:group-theory-applied:3}
$C_{2v}$ में \omterm{def:b3:group-theory-applied:direct-product}{प्रत्यक्ष गुणनफल} $B_1\otimes B_2$, $A_2\otimes B_1$ परिकलित कीजिए, और
$E\otimes E$, $C_{3v}$ में (अंतिम का लघुकरण कीजिए)।
\end{exercise}

\begin{exercise}[$\star$]\label{exo:b3:group-theory-applied:4}
% ledger: vib:CH4.nu1, vib:CH4.nu2, vib:CH4.nu3, vib:CH4.nu4
\ce{CH4} के मूल कंपन 2917 ($a_1$), 1534 ($e$), 3019 और 1306 (दोनों $t_2$)
\unit{cm^{-1}} पर हैं। कौन-से अवरक्त स्पेक्ट्रम में प्रकट होते हैं, कौन-से रामन
स्पेक्ट्रम में?
\end{exercise}

\begin{exercise}[$\star\star$]\label{exo:b3:group-theory-applied:5}
\ce{BF3} के लिए $\Gamma_{\mathrm{vib}} = A_1' + 2E' + A_2''$ को स्थान पर छोड़े गए परमाणुओं से
व्युत्पन्न कीजिए, फिर अवरक्त और रामन बैंड गिनिए। कौन-सी विधा दोनों में से केवल एक
स्पेक्ट्रम में दिखती है, और कौन-सी दोनों में?
\end{exercise}

\begin{exercise}[$\star\star$]\label{exo:b3:group-theory-applied:6}
\ce{SF6} कितने अवरक्त और रामन बैंड दिखाता है? कौन-सी विधा मूक है? कोई
बैंड दोनों स्पेक्ट्रमों में क्यों नहीं दिखता?
\end{exercise}

\begin{exercise}[$\star\star$]\label{exo:b3:group-theory-applied:7}
$h_2$ पर \omterm{def:b3:group-theory-applied:projection}{प्रक्षेप संकारक} लगाइए, $E$ निरूपण के लिए जो $C_{3v}$ का है, और
दिखाइए कि परिणाम का $2h_1 - h_2 - h_3$ के सापेक्ष लांबिकीकरण $h_2 - h_3$ देता है।
\end{exercise}

\begin{exercise}[$\star\star$]\label{exo:b3:group-theory-applied:8}
फ़ैक- और मेर-\ce{ML3(CO)3} के, तथा
\ce{M(CO)6} के अवरक्त CO बैंडों की संख्या की भविष्यवाणी कीजिए।
\end{exercise}

\begin{exercise}[$\star\star$]\label{exo:b3:group-theory-applied:9}
दिखाइए कि अष्टफलकीय संकुल के लिगैंडों के छह $\sigma$ कक्षक
$A_{1g} + E_g + T_{1u}$ का विस्तार करते हैं, और बताइए कि धातु के कौन-से कक्षक हर एक के साथ संयोजित हो सकते हैं।
\end{exercise}

\begin{exercise}[$\star\star\star$]\label{exo:b3:group-theory-applied:10}
नियमित निरूपण के लिए $\chi(E) = h$ और $\chi(R) = 0$, जब $R \ne E$। \omterm{thm:b3:group-theory-applied:reduction}{लघुकरण सूत्र} से
दिखाइए कि उसमें हर \omterm{def:b3:point-groups:irrep}{अलघुकरणीय
निरूपण} $\Gamma_i$ ठीक $d_i$ बार आता है, और $\sum_id_i^2 = h$ निकालिए।
\end{exercise}

\begin{exercise}[$\star\star\star$]\label{exo:b3:group-theory-applied:11}
ऐसीटिलीन \ce{HC#CH} रैखिक है। $D_{2h}$ में कार्य करते हुए $\Gamma_{\mathrm{vib}}$ ज्ञात कीजिए, उसे
$D_{\infty h}$ के अंकनों में पुनः समूहित कीजिए, और अवरक्त तथा रामन स्पेक्ट्रमों की भविष्यवाणी कीजिए।
\end{exercise}

\begin{exercise}[$\star\star\star$]\label{exo:b3:group-theory-applied:12}
जल ($xz$ तल) में इनमें से कौन-से एक-इलेक्ट्रॉन उन्नयन अनुमत
संक्रमण देते हैं, और किस ध्रुवण के साथ: $1b_2 \to 4a_1$, $3a_1 \to 4a_1$,
$1b_2 \to 2b_1$, $1b_1 \to 4a_1$?
\end{exercise}

\section{समस्या: सिस या ट्रांस? स्पेक्ट्रम से पढ़ा गया एक कार्बोनिल संकुल}

\begin{problem}[{सप्ताहांत समस्या --- चार समावयवियों के बिंदु समूह और CO-तनन
निरूपण, उनके अवरक्त और रामन बैंड, और मापे गए स्पेक्ट्रम से
एक उत्पाद की पहचान}]
\label{pb:b3:group-theory-applied:1}
% ledger: ct:C2v, ct:C3v, ct:D4h
एक अष्टफलकीय धातु M कार्बोनिल लिगैंड और अन्य लिगैंड L वहन करती है (हर L को
एक बिंदु के रूप में गिना जाता है)। चार समावयवियों पर विचार किया जाता है: cis- और
trans-\ce{ML4(CO)2}, तथा फ़ैक- और मेर-\ce{ML3(CO)3}। मान लीजिए कि L लिगैंड
सममिति को उससे आगे नहीं घटाते जितना उनकी स्थितियाँ आरोपित करती हैं।
\ce{ML3(CO)3} का एक संश्लेषण ऐसा उत्पाद देता है जिसका अवरक्त स्पेक्ट्रम 2048, 1965 और
\qty{1938}{cm^{-1}} पर तीन प्रबल बैंड दिखाता है (समस्या के आँकड़े); उसका रामन
स्पेक्ट्रम लगभग उन्हीं स्थितियों पर तीन बैंड दिखाता है।

\medskip\noindent\textbf{भाग I --- \omterm{def:b3:point-groups:point-group}{बिंदु समूह}।}
\begin{enumerate}
  \item चारों समावयवियों को एक अष्टफलक पर खींचिए।
  \item cis-\ce{ML4(CO)2} का \omterm{def:b3:point-groups:point-group}{बिंदु समूह} दीजिए।
  \item trans-\ce{ML4(CO)2} का।
  \item फ़ैक-\ce{ML3(CO)3} का।
  \item मेर-\ce{ML3(CO)3} का।
  \item चारों में से किनमें \omterm{def:b3:point-groups:inversion}{प्रतिलोमन केंद्र} है?
\end{enumerate}

\medskip\noindent\textbf{भाग II --- CO तनन।}
\begin{enumerate}[resume]
  \item समझाइए कि C--O आबंध सदिश ऐसा निरूपण क्यों वहन करते हैं जिसमें कोई
        संक्रिया स्थान पर छोड़े गए हर CO के लिए 1 का योगदान देती है।
  \item सिस समावयवी का $\Gamma_{\mathrm{CO}}$ ज्ञात कीजिए।
  \item ट्रांस समावयवी का।
  \item फ़ैक समावयवी का।
  \item मेर समावयवी का।
  \item हर विमा को CO लिगैंडों की संख्या से जाँचिए।
\end{enumerate}

\medskip\noindent\textbf{भाग III --- सक्रियता।}
\begin{enumerate}[resume]
  \item हर समावयवी के \omterm{def:b3:group-theory-applied:activity}{अवरक्त-सक्रिय} CO तनन गिनिए।
  \item \omterm{def:b3:group-theory-applied:activity}{रामन-सक्रिय} तनन गिनिए।
  \item कौन-सा समावयवी \omterm{prop:b3:group-theory-applied:mutual-exclusion}{परस्पर अपवर्जन नियम} दिखाता है?
  \item ट्रांस समावयवी में \omterm{def:b3:group-theory-applied:activity}{अवरक्त-सक्रिय} विधा का वर्णन कीजिए: क्या दोनों CO
        समान कला में तनित होते हैं या विपरीत कला में?
  \item फ़ैक समावयवी में तीन CO केवल दो बैंड क्यों देते हैं?
  \item 1 (अवरक्त) की बैंड-गिनती आपको क्या बताती?
\end{enumerate}

\medskip\noindent\textbf{भाग IV --- उत्पाद।}
\begin{enumerate}[resume]
  \item उत्पाद का अवरक्त स्पेक्ट्रम किस समावयवी की ओर संकेत करता है?
  \item क्या रामन स्पेक्ट्रम संगत है?
  \item सबसे ऊँचा बैंड CO समूहों का समान-कला तनन है। उस समावयवी में यह
        किस \omterm{def:b3:point-groups:irrep}{अलघुकरणीय निरूपण} का है?
  \item क्या स्पेक्ट्रम दोनों समावयवियों के मिश्रण से आ सकता था? कौन-सा
        अतिरिक्त मापन निर्णय करेगा?
  \item परिणाम बताइए: मेर समावयवी के \omterm{def:b3:group-theory-applied:activity}{अवरक्त-सक्रिय} CO तननों की संख्या, फ़ैक
        समावयवी की संख्या के विरुद्ध।
\end{enumerate}
\end{problem}
